% Generated by tools.wolfram_oracle -- do not edit.
\begin{table}[H]\centering\small
\caption{The model document's claims, tested}\label{tab:bm-claims}
\begin{adjustbox}{max width=\textwidth}
\begin{tabular}{p{20mm}p{42mm}p{74mm}l}
\toprule
Claim & Stated & Benchmark finds & Outcome \\
\midrule
\S7.1(1) & lag-1 autocorrelation $\approx0.85$ (non-tidal) & 0.981 without floods; the model gives 0.967 and 0.916: the seasonal cycle dominates & \textbf{not met} \\
\S7.1(3) & days above FloodAlert $\approx 6\max(n_3,1)$ & 16.3 and 129.7 for $n_3=0, 8$, against 6 and 48 & \textbf{not met} \\
\S8.4 & single-day flood spikes & days above FloodAlert per flood: 2.70, 1.75, 1.27 (non-tidal, partial, tidal) & \textbf{not met} \\
\S5.6 & January, July means $\approx\pm A_s$ & 0.985--0.988 of $A_s$ & met \\
\S4.5 & reproduces the observed severe frequency & $n_3$ counts 5 years but is spread over 50: 8 observed per 5 years become 0.8 & \textbf{not met} \\
\S5.4, \S7.3 & byte-identical regeneration & one gauge id seeds 4181464545 in one process and 3341949591 in another; the series also ends at the run date & \textbf{not met} \\
\S6.1 & lag-7 autocorrelations 0.32, 0.028 & 0.3206, 0.0280 & met \\
\S7.2 & mean levels below FloodAlert & at most 1.245\,m & met \\
\S3.5 & a 14.8-day spring--neap cycle & its phase is the day of the year, so it jumps 9.8 days each 1 January & finding \\
\bottomrule
\end{tabular}
\end{adjustbox}
\end{table}
